%-----------------------------------------------------------------------
\subsection{Instance 4: Concept Probe Impossibility}
\label{sec:instance-concept}
%-----------------------------------------------------------------------

\paragraph{Motivation.}
Concept-based explanations---most prominently Testing with Concept Activation
Vectors (TCAV; \citealt{kim2018tcav})---attempt to explain a neural network's
behavior in human-interpretable terms by identifying \emph{concept directions}
in representation space.  For example, the direction corresponding to
``stripes'' in an image classifier's penultimate layer is found by fitting a
linear probe that separates activations on striped versus non-striped inputs.
This direction is then interpreted as the network's internal encoding of the
concept, and the degree to which prediction scores change along it is reported
as the explanation.

The hope is that concept directions are a stable, faithful window into model
internals.  We show this hope is mathematically impossible under
underspecification: two models with identical predictions can encode the same
concept in directions that are genuinely incompatible, so no concept probe can
simultaneously be faithful, stable, and decisive.

\paragraph{Instantiation.}
We instantiate Definition~\ref{def:explanation-system} as follows.

\begin{definition}[Concept Probe Explanation System]
\label{def:concept-system}
Let $c \in \mathbb{N}$ be the dimension of the representation space.  The
\emph{concept probe explanation system} is the tuple
$\mathcal{S}_{\mathrm{cp}} = (\Params, \sH, Y, \mathsf{obs}, \mathsf{exp}, \incomp)$ where:
\begin{itemize}
  \item $\Params = \Theta$ is the space of model weight configurations,
    each $\theta \in \Theta$ specifying a neural network together with its
    internal representations;
  \item $\sH = \R^c$ is the space of \emph{concept activation vectors}---unit
    directions in the $c$-dimensional representation space that a linear probe
    associates with a target concept;
  \item $Y$ is the space of predictions (input-output functions
    $f_\theta : \mathcal{X} \to \mathcal{Y}$), serving as the observable
    behavior of the model;
  \item $\mathsf{obs}(\theta) = f_\theta$ is the \emph{observation map},
    recording the input-output function computed by weights $\theta$;
  \item $\mathsf{exp}(\theta) = v_\theta \in \R^c$ is the \emph{explanation
    map}, returning the concept activation direction extracted by a linear
    probe on the representations of $\theta$; and
  \item $(v, v') \in \incomp$ iff $v$ and $v'$ are not related by a rotation
    within the concept subspace---that is, they designate structurally distinct
    directions that cannot both be reported as \emph{the} concept direction
    without contradiction.
\end{itemize}
\end{definition}

\paragraph{The Rashomon Property for concept probes.}
The empirical literature provides direct evidence that the Rashomon property
holds for concept directions.  \citet{bolukbasi2016man} demonstrate that gender
directions vary substantially across word-embedding architectures trained on the
same corpus: models with identical downstream task performance encode gender in
directions that differ in orientation, dimension, and linearity.
\citet{kim2018tcav} note that TCAV concept directions are representation-dependent
and can shift when the network is retrained with a different random seed.
\citet{ravfogel2020null} show via iterative nullspace projection that a single
concept is not localized to one direction: removing the dominant linear probe
direction still leaves residual concept structure, and the secondary directions
are mutually incompatible with the primary.

Formally, we require:
\begin{definition}[Rashomon Property for Concept Probes]
There exist weight configurations $\theta_1, \theta_2 \in \Theta$ such that
\[
  f_{\theta_1} = f_{\theta_2}
  \quad \text{(same predictions)}
  \qquad \text{and} \qquad
  (v_{\theta_1},\, v_{\theta_2}) \in \incomp
  \quad \text{(incompatible concept directions).}
\]
\end{definition}

This is the formal content of Definition~\ref{def:rashomon} in
$\mathcal{S}_{\mathrm{cp}}$: equivalent models (same observable behavior,
hence $\mathsf{obs}(\theta_1) = \mathsf{obs}(\theta_2)$) report incompatible
concept directions.  The condition holds in any overparameterized network
class---the zero-loss manifold has dimension $\geq d - n > 0$ for hidden
dimension $d > n$ training points, so there are infinitely many weight
configurations with identical input-output behavior but geometrically distinct
internal representations.

\paragraph{The concept probe impossibility.}

\begin{theorem}[Concept Probe Impossibility]
\label{thm:concept-impossibility}
If the concept probe explanation system $\mathcal{S}_{\mathrm{cp}}$ has the
Rashomon property, then no explanation map $\mathsf{exp}$ can simultaneously
be faithful, stable, and decisive.
\end{theorem}

\begin{proof}
Immediate from Theorem~\ref{thm:universal} applied to
$\mathcal{S}_{\mathrm{cp}}$.  Concretely: let $\theta_1, \theta_2$ witness the
Rashomon property, so $\mathsf{obs}(\theta_1) = \mathsf{obs}(\theta_2)$ and
$(v_{\theta_1}, v_{\theta_2}) \in \incomp$.  Faithfulness forces
$\mathsf{exp}(\theta_i) = v_{\theta_i}$; decisiveness requires each
$\mathsf{exp}(\theta_i)$ to be a specific direction; stability forces
$\mathsf{exp}(\theta_1) = \mathsf{exp}(\theta_2)$.  But then
$v_{\theta_1} = v_{\theta_2}$, contradicting $(v_{\theta_1}, v_{\theta_2}) \in
\incomp$.
\end{proof}

The Lean~4 formalization of this instance is in
\texttt{UniversalImpossibility/ConceptInstance.lean}.  The theorem
\texttt{concept\_impossibility} is proved by a single application of
\texttt{explanation\_impossibility} with zero additional axioms beyond the
Rashomon property, which is stated as the axiom \texttt{conceptSystem} encoding
the specific system tuple.

\paragraph{Implications for probing-as-understanding.}
Theorem~\ref{thm:concept-impossibility} has direct consequences for the
practice of using concept probes as explanations of model internals.

\begin{enumerate}
\item \emph{Concept directions are not model-invariant facts.}
  A TCAV direction is a fact about a specific weight configuration, not about
  the function the network computes.  Two retraining runs that produce
  functionally identical models (same accuracy, same predictions on held-out
  data) may produce concept directions that are not related by rotation and
  therefore cannot be reconciled into a single explanation.

\item \emph{Faithfulness and stability cannot coexist.}
  Requiring a probe to faithfully report the concept direction of the model it
  was fitted to forces instability across retraining seeds.  Requiring
  stability forces the probe to report a direction that is unfaithful to at
  least one model in the Rashomon set.

\item \emph{The resolution entails indecisiveness.}
  The $G$-invariant resolution (Section~\ref{sec:resolution}) applied here
  would average concept directions across equivalent models, yielding a
  distribution over directions or a subspace rather than a single vector.
  This restores faithfulness-in-expectation and stability at the cost of
  decisiveness: the explanation no longer commits to a single direction, which
  undermines the interpretive clarity that motivates concept probing in the
  first place.
\end{enumerate}

The impossibility does not mean concept probes are uninformative, but it does
mean that treating a single concept direction as \emph{the} faithful,
reproducible encoding of a concept in a model is unjustified under
underspecification.  Practitioners should report concept directions together
with their variability across the Rashomon set---analogously to how
\DASH{} reports SHAP values with their ensemble distribution rather than as a
single point estimate.
